\documentclass[sigconf,nonacm]{acmart}
\settopmatter{printacmref=false, printccs=false, printfolios=true}

\usepackage{booktabs}
\usepackage{multirow}
\usepackage{xurl}

% Sigconf is a narrow two-column layout; long identifier-like tokens
% inside \texttt / \textsf (slashes, hyphens, dots) have no natural
% breakpoints and cause column overflow. \sloppy plus a large
% emergency stretch lets TeX find acceptable line breaks; the macros
% below insert discretionary break points after the typical
% identifier separators for places where the long token still does
% not fit (e.g.\ a 50-character file path).
\sloppy
\setlength{\emergencystretch}{3em}
\newcommand*{\Bsl}{\discretionary{/}{}{/}}% breakable slash
\newcommand*{\Bhy}{\discretionary{-}{}{-}}% breakable hyphen

% Hide the ACM submission boilerplate that has no place in a
% self-published technical report.
\setcopyright{none}
\acmConference{}{}{}
\renewcommand{\footnotetextcopyrightpermission}[1]{}
\pagestyle{plain}

\title{gheim-ch-pii-212k and gheim-ch-560m: An Open Swiss-Language PII Dataset and NER Model. Technical report.}
\author{Joel Barmettler}
\affiliation{%
  \institution{Independent}
  \city{Zürich}
  \country{Switzerland}
}
\email{jbarmettler@proton.me}

\begin{document}

\begin{abstract}
We release \textsf{gheim-ch-pii-212k}, a 212{,}503-chunk personally-identifiable
information (PII) named-entity-recognition dataset covering the four
official Swiss languages (Swiss German, Swiss French, Swiss Italian,
Romansh) and English, and \textsf{gheim-ch-560m}, a 560M-parameter
\textsf{xlm-roberta-large} fine-tune released alongside it. The dataset
combines real Swiss text from court rulings, federal parliament records,
Swiss-filtered web text and a Romansh corpus with a templated and
LLM-generated synthetic gap-fill layer covering cells where the
real-text corpus was too sparse for training and evaluation.
Annotations come from a three-LLM consensus (Gemma~4 26B-A4B,
Qwen3.6 35B-A3B, Nemotron-3 Nano Omni 30B-A3B) plus a regex catalogue
with checksum validators, merged with per-span provenance.
\textsf{gheim-ch-560m} attains 0.70--0.94 person-NER char F\textsubscript{1}
on four external benchmarks
(\textsf{ai4privacy/openpii-1m}, ZurichNLP \textsf{swissner},
Babelscape \textsf{WikiNeural}, CoNLL-2003) without ever seeing
training data from any of them; on the in-distribution held-out test
split it attains a strict-span F\textsubscript{1} of 0.910 (char
F\textsubscript{1} 0.946). Among open Swiss-applicable detectors run
zero-shot on the same Swiss test set, the strongest comparable
overall-F\textsubscript{1} entry is Microsoft Presidio at 0.398
strict / 0.501 char.
The dataset ships under CC BY 4.0; the model under Apache 2.0; both
are available on the Hugging Face Hub. Reproducible eval scripts and
checkpoints are on GitHub. This is an engineering report on a
publicly-released artifact; no new method or architecture is proposed.
\end{abstract}

\maketitle

% --------------------------------------------------------------------
\section{Introduction}
\label{sec:intro}

The dominant pattern in production large-language-model (LLM) workflows
is to call a hosted API: OpenAI, Anthropic, Google, Mistral. For
Swiss-market deployments this raises a recurring compliance concern:
every request that contains a customer name, a Swiss address, an IBAN,
or an AHV/AVS social-security number leaves Swiss jurisdiction the
moment it crosses the API boundary. The legal landscape for using
third-party LLM services with personal data under Swiss law is
analysed in detail by~\citet{rosenthal2025jusletter}; the practical
mitigation in production systems is to anonymise sensitive content
before the request and restore the originals on the response. The
accuracy of the round-trip depends on the quality of the underlying
named-entity recogniser.

Existing public PII detectors handle Swiss text poorly. The closest
open generalist, \textsf{openai/privacy-filter}~\cite{openai2025privacyfilter}
(an OLMoE-based~\cite{muennighoff2024olmoe} 1.4B-parameter
mixture-of-experts model with 50M active parameters), reaches a
strict-span F\textsubscript{1} of 0.44 on our held-out Swiss test
set despite multilingual training. Microsoft Presidio~\cite{microsoft2023presidio},
configured with per-language NLP engines, reaches 0.43. Multilingual
NER baselines such as Davlan's \textsf{xlm-roberta-base-ner-hrl}
~\cite{davlan2022xlmrnerhrl} reach 0.65 person-only F\textsubscript{1}.
None of the available systems fluently handle the four official
Swiss languages together with the structured PII formats common in
Swiss documents (CH-IBAN, AHV, VAT-CHE), and only one
(SwissNER~\cite{vamvas2023swissner}) provides a Romansh evaluation
slice.

This report documents three publicly-released artifacts:

\begin{enumerate}
  \item \textsf{gheim-ch-pii-212k}, a multilingual PII NER dataset of
        212,503 chunks across Swiss German, Swiss French, Swiss Italian,
        Romansh and English, covering eight PII categories
        (account number, address, date, email, person, phone, URL,
        secret), encoded in a 33-class BIOES tag scheme byte-aligned
        with \textsf{openai/privacy-filter}.
  \item \textsf{gheim-ch-560m}, a 560M-parameter token-classification
        model fine-tuned from \textsf{xlm-roberta-large}~\cite{conneau2020unsupervised}
        on the dataset, attaining a held-out test F\textsubscript{1}
        of 0.910 strict-span and 0.946 character-level.
  \item A reproducible evaluation harness that scores the released
        model on four established external benchmarks
        (\textsf{ai4privacy\Bsl openpii\Bhy 1m}, ZurichNLP \textsf{swissner},
        Babelscape \textsf{WikiNeural}~\cite{tedeschi2021wikineural},
        CoNLL\Bhy 2003~\cite{tjong2003introduction}) and seven competing
        open detectors on the released test set, with replication
        checks against each baseline's published numbers on its native
        dataset.
\end{enumerate}

\subsection{Scope}
\label{sec:scope}

This is a technical report on a deployed artifact, not a research
contribution. We do not propose new model architectures, training
methods, evaluation metrics, or annotation schemes. The model is
\textsf{xlm-roberta-large}~\cite{conneau2020unsupervised} fine-tuned
with stock AdamW~\cite{loshchilov2019decoupled} and a small
hyperparameter sweep; the label space is byte-identical to
\textsf{openai/privacy-filter}~\cite{openai2025privacyfilter}; the
real-text corpora are taken as-is from the Apertus pretrain
release~\cite{apertus2025} and downstream sources; the dataset
construction pattern (LLM annotation followed by checksum-validated
regex augmentation and template-based gap-fill) is by now standard
practice. The contribution is the assembled artifact: a
publicly-released Swiss-language PII dataset that includes Romansh,
a strong out-of-the-box checkpoint for the same niche, and a
reproducible evaluation that puts both into context against the open
alternatives. The report documents what the dataset contains, how
the model was trained, what its measured behaviour is on multiple
test sets, and what its known limitations are, so downstream users
can judge whether to deploy it.

The rest of the paper is structured as follows.
Section~\ref{sec:related} situates the work against prior art.
Section~\ref{sec:dataset} describes the dataset construction.
Section~\ref{sec:model} covers model selection and training.
Section~\ref{sec:eval} reports the evaluation.
Section~\ref{sec:limits} discusses limitations and
Section~\ref{sec:conclusion} summarises.
Appendix~\ref{app:results} collects the four full result tables
(per-language $\times$ per-category breakdown, comparison to other
detectors, cross-domain transfer, and a methodology-validation pass
on each external baseline).

% --------------------------------------------------------------------
\section{Related work}
\label{sec:related}

\subsection{PII detection and de-identification}

Microsoft Presidio~\cite{microsoft2023presidio} is a widely-used
industrial open-source PII detector: a per-language spaCy NER
engine combined with checksum-validated regex recognisers for
structured formats (IBAN, credit card, SSN, etc.). Its built-in
catalogue is large and the framework is easy to extend; its weakness
on Swiss text is inherited from the per-language spaCy NER
backbone (Section~\ref{sec:eval}, Table~\ref{tab:directionA}).

The AI4Privacy family of corpora and fine-tuned
checkpoints~\cite{ai4privacy2024openpii,ai4privacy2024pii200k} is a
high-volume recent open-source effort. Their datasets are templated
synthetic prose with embedded Faker-generated PII at the million-chunk
scale across 12+ fine-grained categories and 17+ locale variants,
and the published distilbert fine-tunes (e.g.\ \textsf{Isotonic\Bsl distil\Bhy bert\_\allowbreak finetuned\_\allowbreak ai4privacy\_\allowbreak v2})
are widely-downloaded open PII detectors on the Hugging Face Hub.
Section~\ref{sec:eval} confirms the published per-token accuracy
numbers but shows a sharp drop on real Swiss text: when scored under
strict-span F\textsubscript{1} on \textsf{gheim-ch-pii-212k} the
fine-tunes reach 0.16. The cross-domain gap reflects template-vs-real
domain shift more than label-noise: AI4Privacy's underlying templates
were not designed with Swiss surface forms (CH-IBAN, AHV/AVS,
\texttt{+41} phones, German-language addresses) as targets.

\textsf{openai/privacy-filter}~\cite{openai2025privacyfilter} is the
strongest open zero-shot generalist on our test set. It is a
token-classification head atop OLMoE~\cite{muennighoff2024olmoe},
a 1.4B-total-parameter mixture-of-experts encoder with 50M parameters
active per token, and it shares our 33-class BIOES schema. Its
\texttt{openai}-namespace release made it the default upstream from
which our schema is byte-aligned, so a fine-tune started from this
model can reuse the pretrained classifier-head weights.

\subsection{Multilingual NER}

\textsf{XLM-RoBERTa}~\cite{conneau2020unsupervised} is a common
starting point for multilingual token-classification and is the base
of \textsf{gheim-ch-560m}. The Davlan family of
fine-tunes~\cite{davlan2022xlmrnerhrl} are trained on PER/ORG/LOC
across many European languages and provide our generalist baseline.
WikiNeural~\cite{tedeschi2021wikineural} provides large-scale
silver-data multilingual NER training and evaluation across nine
European languages and is one of our cross-domain external
benchmarks. The CoNLL-2003 shared
task~\cite{tjong2003introduction} is a long-running English NER
benchmark and is also part of our cross-domain evaluation. The BIOES
tag scheme used throughout follows the design analysis of
\citet{ratinov2009design}.

\subsection{Swiss-NLP and the Apertus stack}

\textsf{SwissBERT}~\cite{vamvas2023swissbert} is an encoder
pretrained on Swiss text covering all four official languages,
including Romansh. We considered SwissBERT as a base in the bake-off
(Section~\ref{sec:model}, Table~\ref{tab:bakeoff}) and selected
XLM-R-large by 0.8 pp validation F\textsubscript{1}; SwissBERT is a
competitive alternative when ${\sim}$300M parameters is the
deployment ceiling. Its companion test
set, \textsf{SwissNER}~\cite{vamvas2023swissner}, is the only public
Swiss-news NER set that includes Romansh and is one of our
cross-domain external benchmarks. The Romansh evaluation in particular
relates closely to the \textsf{WMT24++} benchmark
extension of~\citet{vamvas2025romansh_wmt}, which adds Rumantsch
Grischun together with the five regional Romansh idioms (Sursilvan,
Sutsilvan, Surmiran, Puter, Vallader) to the WMT machine-translation
benchmark; the Apertus Romansh corpus, which underlies our
\texttt{rm} chunks, is constructed in a similar idiom-aware way.

The real-text portion of our dataset is sourced from the Apertus
pretrain corpora released alongside the Apertus Swiss-language
LLM~\cite{apertus2025}. Apertus aggregates four Swiss-domain sources:
the Open Legal Data Platform \textsf{Entscheidsuche}, which provides
machine-readable Swiss court decisions across cantons in their
original German, French, or Italian~\cite{entscheidsuche2025};
\textsf{Curia Vista}, the official business database of the Swiss
Federal Assembly, which includes parliamentary motions, postulates,
and transcripts in German, French, and Italian~\cite{curiavista2024};
the Swiss-filtered subset of FineWeb-2~\cite{penedo2025fineweb2} (the
multilingual web-crawl corpus that scales the FineWeb pipeline of
\citet{penedo2024fineweb} to all languages and applies XLM-R-based
quality filters); and a purpose-built Romansh corpus drawing on
cantonal law texts, the Lia Rumantscha online dictionaries, the
ZurichNLP bilingual corpus, and Romansh Wikipedia. These corpora
are released under CC-BY-4.0-compatible licences and are the
substrate for the real-Swiss-text portion of \textsf{gheim-ch-pii-212k}.

\subsection{Quantisation for browser deployment}

The released model ships an int8 ONNX export so the demo at
\href{https://gheim.ch}{gheim.ch} runs entirely in-browser via
WebGPU or WebAssembly. The quantisation procedure follows the
post-training integer-arithmetic-only inference framework of
\citet{jacob2018quantization}, applied through ONNX
Runtime~\cite{onnxruntime} via the \textsf{optimum} dynamic-quantizer
with AVX-512 VNNI as the deployment target. The browser runtime is
\textsf{Transformers.js}~\cite{transformersjs}.

\subsection{Industry tooling}

\textsf{transformers}~\cite{wolf2020transformers} is the training
and inference runtime. \textsf{seqeval}~\cite{nakayama2018seqeval}
provides the strict-span scoring used both for our headline
F\textsubscript{1} and for every comparison and replication entry
in Section~\ref{sec:eval}. \textsf{spaCy}~\cite{honnibal2020spacy}
provides the per-language NER pipelines used in our comparison
and inside Presidio. AdamW~\cite{loshchilov2019decoupled} is the
optimiser. Faker\_CH provides the Swiss-locale random data used in
our synthetic gap-fill payload generation.

% --------------------------------------------------------------------
\section{Dataset: gheim-ch-pii-212k}
\label{sec:dataset}

\subsection{Source corpora and language coverage}

Real-text content is drawn from the Apertus pretrain
corpora~\cite{apertus2025}: the Swiss legal-decisions
subset \texttt{entscheidsuche} curated from the Open Legal Data Platform
of the same name~\cite{entscheidsuche2025}, federal-parliament
proceedings drawn from the \texttt{curia\_vista} business
database~\cite{curiavista2024}, the Swiss-filtered slice of
FineWeb-2~\cite{penedo2025fineweb2}, and the Romansh corpus released
alongside Apertus, which compiles material from cantonal law texts,
the Lia Rumantscha online dictionaries, the ZurichNLP bilingual
corpus, and Romansh Wikipedia. Source documents are chunked at
sentence boundaries to roughly 1{,}500 characters, then tagged with
\texttt{facebook/fasttext-language-identifica\-tion} at confidence
$\geq 0.5$. Chunks below threshold are dropped.
The five-language distribution of the released dataset is shown in
Table~\ref{tab:lang-dist}.

\begin{table}[h]
\centering
\caption{Per-language chunk distribution.}
\label{tab:lang-dist}
\small
\begin{tabular}{l l r r}
\toprule
Code   & Variant                       & Chunks & \% \\
\midrule
de\_ch & Swiss German (written std.)   & 90{,}292 & 42.5 \\
fr\_ch & Swiss French                  & 64{,}021 & 30.1 \\
it\_ch & Swiss Italian                 & 36{,}147 & 17.0 \\
rm     & Romansh                       & 16{,}514 & ~7.8 \\
en     & English (Swiss-context)       & ~5{,}529 & ~2.6 \\
\midrule
\multicolumn{2}{l}{Total}              & 212{,}503 & 100 \\
\bottomrule
\end{tabular}
\end{table}

Spoken-form Swiss German (GSW) is not represented: the corpus is
written-standard German as used in Swiss legal and journalistic
contexts. The fasttext language identifier labels GSW
(``chuchichäschtli''-style orthography) as standard German, so
incidental dialect content cannot be reliably separated.

\subsection{Annotation schema}

The label space follows the \textsf{openai/privacy-filter} taxonomy:
eight PII categories listed in Table~\ref{tab:cats}, each encoded in a
BIOES tag scheme~\cite{ratinov2009design} (Begin, Inside, Outside, End,
Single), giving 33 classes overall (one outside class plus four BIOES
positions for each of the eight categories). The byte-identical label
mapping was chosen for compatibility: it lets a fine-tune started
from \textsf{openai/privacy-filter}'s checkpoint reuse the pretrained
classifier-head weights without remapping, and it lets downstream
consumers swap the two checkpoints in place. The cost is that the
upstream taxonomy's choices are inherited as-is, including ones we
would have made differently. \texttt{account\_number} lumps four
structurally distinct formats (CH-IBAN, AHV/AVS, VAT-CHE, credit
cards) into one cell, so the model cannot signal which kind it
detected; \texttt{secret} is broad enough to catch any high-entropy
token that looks key-shaped; there are no \texttt{organization} or
\texttt{location} cells, which some downstream tasks would want. The
\texttt{private\_} prefix is historical; all eight cells are applied
with a recall-oriented labelling policy and may flag publicly-listed
institutional entities (e.g.\ a court switchboard phone number).

\begin{table}[h]
\centering
\caption{The eight PII categories.}
\label{tab:cats}
\small
\begin{tabular}{l l}
\toprule
Category               & Examples \\
\midrule
\texttt{account\_number} & CH-IBAN, AHV/AVS, VAT-CHE, credit cards \\
\texttt{private\_address} & Bahnhofstrasse 1, 8001 Zürich \\
\texttt{private\_date}   & 1990-01-02, 14.\,M\"arz 2026 \\
\texttt{private\_email}  & alice@example.ch \\
\texttt{private\_person} & Joel Barmettler, Dr.\ Schmidt \\
\texttt{private\_phone}  & +41 44 268 12 34 \\
\texttt{private\_url}    & https://kanzlei.ch/...{} \\
\texttt{secret}          & API keys, bearer tokens \\
\bottomrule
\end{tabular}
\end{table}

\subsection{Curation pipeline}

\paragraph{Labelling.} Real-text content is annotated by three
independent open-weights LLM labellers run in parallel on the same
inputs: \textsf{Gemma 4 26B-A4B} (4-bit AWQ via vLLM,
\texttt{cyankiwi/gemma\Bhy 4\Bhy 26B\Bhy A4B\Bhy it\Bhy AWQ\Bhy 4bit}),
\textsf{Qwen3.6 35B-A3B} (4-bit AWQ), and \textsf{Nemotron-3 Nano
Omni 30B-A3B Reasoning} (FP8). Each labeller is constrained to the
same JSON-schema output so the spans are directly parseable. The
first pass used Gemma alone; a small-sample human-eyeball gate
(n=176) showed F\textsubscript{1}=0.66 against subagent-resolved gold,
with the bulk of the gap in commercial-register director listings
and recall-aggressive court-role text. We added Qwen as a second
labeller to lift recall on those failure modes, then Nemotron as a
third for tie-breaking when Gemma and Qwen disagreed. The three
labellers were applied to the full Apertus pretrain corpora
(approximately 4.5 million chunks across the five languages); the
combined wall-clock cost is roughly 30 GPU-hours on a
2$\times$RTX 4090 setup.

\paragraph{Regex augmentation.} A separate pass re-scans the same
Apertus corpora with a regex catalogue gated by checksum validators
(ISO 13616 mod-97 for IBAN, EAN-13 mod-10 for AHV, ISO 7064 mod-11
for VAT-CHE, Luhn for credit cards). This is an independent scan
rather than a post-filter on the LLM output, so it can pick up
structured PII in chunks that the LLMs rejected or missed. A
column-context heuristic gates the regex hits so that numeric
tables (court statistics, parliamentary vote counts) whose digits
happen to be IBAN- or AHV-shaped are not mis-flagged as
\texttt{account\_number}; the audit in
Section~\ref{sec:labelnoise} shows that the heuristic does not
catch every such case, and it is one of the open hardening
directions.

\paragraph{Format-aware commercial-register re-labelling.} A
fourth annotator runs only on the commercial-register subset, with
a prompt tuned to that format. This was added after the audit
showed the general-purpose labellers under-recalled director and
board-member names in long contact-block listings.

\paragraph{Merge with per-span provenance.} The four annotator
streams (three LLM, one regex) are merged at the span level. Each
surviving span carries the tuple of signals that produced it and a
confidence score derived from how many sources agreed. Two
post-processing steps run unconditionally on the merged output: a
Geonames-CH gazetteer demotes spans labelled \texttt{private\_person}
whose surface is a known Swiss place name (e.g.\ ``Chur'', ``Sion'')
to the appropriate \texttt{private\_address} sub-span or drops them
when the place stands alone, and a label-conflict resolver collapses
overlapping spans at the same character offsets, preferring the
higher-confidence label. Both steps were added after we observed
the earlier pipeline mis-tagged municipality names as people and
double-labelled the same characters under two different categories.

\paragraph{Balancing.} The merged set of approximately 2.3M
candidate chunks is passed through a five-phase balancer to limit
document concentration, value memorisation and per-language class
imbalance. The cap numbers were chosen from a frequency-analysis
pass over the unbalanced output. Phase E drops spans whose
provenance-derived confidence is below 0.5 (a span with no signal
agreement is unstable enough that we prefer to discard it).
Phase A caps each source document to 30 chunks. Phase B caps each
(language, category, normalised value) tuple to 30 chunks (60 for
the medium-difficulty Swiss-language person/address/date cells,
100 for any Romansh cell, 200 for English and for synthetic
sources), so a single recurring CH-IBAN or person name cannot
dominate its cell. Phase C applies per-(language $\times$ source
$\times$ category) cell caps tiered by data abundance: up to
24{,}000 chunks for the data-rich Swiss-language person, address
and date cells in legal and parliamentary text; 18{,}000 for the
other major cells; 12{,}000 for Italian; 10{,}000 to keep all of
English; smaller caps where the corpus simply does not have the
material. \texttt{secret} is effectively uncapped because it is
rare corpus-wide. Phase D injects negative chunks (chunks with no
detected PII) at 10\,\% of positives per language so the model
trains on enough out-of-class context. The five phases together
bring the candidate set down to the released chunk count.

\paragraph{Synthetic gap-fill.} A separate synthetic layer is
appended to populate cells where real-Swiss-text coverage was
insufficient. Eight templated layers are generated from
hand-written fragment files (greetings, signatures, bodies,
legal/medical/bank/HR documents, multilingual KYC/HR/bank/
secret-leak/doctor-note documents in DE/FR/IT, short-form
narrative in all five languages, DE structured layouts
(CSV/vCard/YAML), common-word surname positive/negative pairs,
adversarial-negative chunks, and a Romansh-context secret-slot
layer). A composer framework substitutes Faker\_CH-generated names,
addresses and phones; checksum-validated CH-IBAN, AHV and VAT-CHE;
and common secret formats into the templates, then optionally
applies one of seven realistic noise overlays (NBSP, broken caps,
no-space, markdown, HTML entity, stray punctuation, OCR swap) at
15\,\% probability per chunk. A slot-fill verifier additionally
prompts Gemma to embed pre-generated payload values verbatim into
natural multilingual prose for the rarer cells; it accepts a chunk
only if every payload value appears as an exact substring in the
model output. Synthetic chunks carry a \texttt{synthetic} boolean
so consumers can filter to real-only or synthetic-only as
required.

\subsection{Splits and held-out integrity}

Documents are partitioned at the source-document level before any
chunk is taken, with stratification by (language $\times$ source), so
no chunk in test or validation comes from a document with a chunk
in train. Synthetic chunks are placed deterministically (random seed
20260522). Final split sizes are 170{,}001 / 21{,}256 / 21{,}246
chunks for train / validation / test.

\subsection{Label-noise audit}
\label{sec:labelnoise}

The dataset's labels are machine-generated. No third-party human
annotators were contracted, and no subset carries human gold labels.
Every model F\textsubscript{1} reported in Section~\ref{sec:eval} is
therefore a measure of agreement between the trained model and its
machine-generated labelling source, not a measure of agreement with
a human ground truth.

To put a number on the residual disagreement we ran a four-way
re-labelling audit. A stratified 580-chunk sample (drawn from the
val and test splits, not the training pool) was re-annotated by
four named 2026-era LLMs accessed via OpenRouter ---
\textsf{moonshotai/kimi\Bhy k2.6}, \textsf{deepseek/deepseek\Bhy v4\Bhy pro},
\textsf{minimax/minimax\Bhy m2.7} and \textsf{z\Bhy ai/glm\Bhy 5.1} ---
each given the same per-language prompt the dataset's own labeller
uses, with \texttt{temperature=0} and \texttt{seed=0} so the run is
reproducible from a checkout. The four per-model labellings are
combined into a per-span majority vote (a span is in the consensus
iff ${\geq}2$ of 4 models agree on it). The released dataset's
labels reach an F\textsubscript{1} of \textbf{0.71} against that
majority vote (precision 0.747, recall 0.680, $n=580$ chunks,
1{,}349 gold spans). Per-category breakdown shows the expected
spread: structured categories where the surface form carries the
signal score high (\texttt{private\_phone} 0.87,
\texttt{private\_email} 0.86, \texttt{private\_url} 0.78),
free-text categories with policy room score lower
(\texttt{private\_date} 0.72, \texttt{private\_person} 0.70),
and the categories whose definition depends most on the
recall-oriented labelling policy score lowest
(\texttt{account\_number} 0.52, \texttt{private\_address} 0.43;
\texttt{secret} 0.67 on only 4 spans). Inspection of the
\texttt{private\_person} disagreements shows the gap is dominated
by policy interpretation rather than random annotator error: the
audit models tend to omit publicly-listed officials, authors of
court rulings, and witness names, while the dataset's
recall-oriented policy includes them.

For comparison, the four audit models also disagree among themselves.
Per-model F\textsubscript{1} against the dataset's labels (no
voting) ranges from 0.42 (\textsf{kimi\Bhy k2.6}, the most
conservative) to 0.73 (\textsf{deepseek\Bhy v4\Bhy pro}, the most
policy-aligned). This range bounds what any single audit run
could plausibly produce, and motivates the majority-vote consensus
as the more stable reference. The raw per-model labellings, the
consensus, and the scoring script are committed under
\texttt{data/audit\_openrouter\_*.jsonl} so the 0.71 number is
reproducible from the released sample.

The 0.71 audit number is best read as a measure of \emph{how
aggressively the dataset labels relative to a four-model
consensus}, not as a label-noise floor in the random-error sense;
the model's headline 0.910 strict F\textsubscript{1} should be read
against this with that caveat. An earlier development pilot
(n=176, single-model Qwen re-annotation) gave F\textsubscript{1}=0.66
and was used during dataset construction as a coarse quality gate;
the 0.71 number above is the reproducible, four-model successor.

\paragraph{What drives the gap.} A natural worry is that the 0.71
F\textsubscript{1} is largely span-boundary noise (commas, titles,
the ``Dr.\ Müller'' vs ``Müller'' fragmentation that NER scoring is
notoriously sensitive to). It is not. Re-scoring with a fuzzy
match that treats any same-label substring overlap as a match
yields F\textsubscript{1}=0.73 --- only \textbf{1.7 pp} of the gap
is boundary noise. The remaining 28 pp falls into three concrete
patterns, each pointing at a specific dataset-quality follow-up
for a future revision:
\begin{itemize}
  \item \textbf{Regex-layer false positives in
        \texttt{account\_number}} (52 dataset-only spans). The
        regex-mining pass picked up numeric tables in court
        documents and parliamentary statistics --- e.g.\
        ``2011 2011 2011 1-3 1-3 1'', ``1892 101 14123 1519'' ---
        as IBAN/AHV-shaped, but they are not PII. All four audit
        models reject them. \emph{This is a residual dataset bug a
        follow-up cleanup pass should tighten further} (the released
        artifact already applies a column-context heuristic; the
        audit shows it does not catch every case).
  \item \textbf{Recall-aggressive labelling of court roles}
        (101 \texttt{private\_person} dataset-only,
        104 \texttt{private\_address} dataset-only). The dataset
        labels witnesses, judges, court reporters, and named
        places-of-origin in legal opinions; the audit models
        follow the prompt's ``do not tag generic roles'' rule
        more strictly. This is a real policy choice the dataset
        makes by design, and a future revision could expose it as
        an opt-in flag rather than the default.
  \item \textbf{Under-recall of names in commercial-register
        text} (156 \texttt{private\_person} consensus-only). The
        audit models consistently flag managing directors and
        board members in Swiss commercial-register listings
        (e.g.\ ``Reto Rubeli'', ``Fabienne Rubeli-Riedo'',
        ``Karim Lassueur'') that the dataset's labelling pass
        missed in long contact-block listings. \emph{This is a
        residual recall gap a follow-up revision could address} by
        re-labelling the commercial-register subset with a prompt
        that emphasises this format (the released artifact already
        runs a format-aware re-labeller on this subset; the audit
        shows further yield is still on the table).
\end{itemize}
The artefacts that produced these numbers (raw per-model
labellings, consensus, classification breakdown, example chunks)
are committed under \texttt{data/audit\_openrouter\_*.jsonl},
\texttt{data/audit\_openrouter\_report.json} and
\texttt{data/audit\_openrouter\_analysis.json}, and the scripts
that produced them are at \texttt{training/src/gheim\_training/}
\texttt{data/analysis/audit\_openrouter.py} and
\texttt{audit\_score.py}.

% --------------------------------------------------------------------
\section{Model: gheim-ch-560m}
\label{sec:model}

\subsection{Architecture and base-model selection}

The model is a token-classification head atop
\textsf{xlm-roberta-large}~\cite{conneau2020unsupervised} (24 layers,
1024 hidden, 16 heads, 560M parameters). XLM-R-large was selected
following a controlled bake-off against \textsf{ZurichNLP/swissbert}
~\cite{vamvas2023swissbert} (270M dense), summarised in
Table~\ref{tab:bakeoff}. Both bases received an identical 5$\times$3
hyperparameter sweep over (learning rate, layer-wise LR decay) at
one epoch, after which the winning configuration per base was trained
for three full epochs and selected by best validation F\textsubscript{1}.
A third candidate, \textsf{openai/privacy-filter}, was deferred:
the released model is required to ship as an ONNX export so that
the in-browser demo at \href{https://gheim.ch}{gheim.ch} can run
via \textsf{transformers.js}, and the Olmo-MoE architecture is not
currently stable in \textsf{optimum}'s~\cite{onnxruntime} ONNX
export path.

\begin{table}[h]
\centering
\caption{Base-model bake-off. Numbers are from the original base-model
selection on the predecessor dataset (\textsf{gheim-ch-pii-171k}); the
released \textsf{xlm-roberta-large} fine-tune on
\textsf{gheim-ch-pii-212k} attains 0.910 val / 0.910 test
strict-span F\textsubscript{1}. The bake-off has not been re-run on
the new dataset because the architecture decision (XLM-R-large over
SwissBERT) was settled at the time of selection.}
\label{tab:bakeoff}
\small
\begin{tabular}{l r r}
\toprule
Base model                                  & Val F\textsubscript{1} & Test F\textsubscript{1} \\
\midrule
\textsf{ZurichNLP/swissbert} (270M)         & 0.910 & 0.910 \\
\textsf{FacebookAI/xlm-roberta-large} (560M)& \textbf{0.918} & \textbf{0.916} \\
\bottomrule
\end{tabular}
\end{table}

In the (LR, LLRD) sweep, learning rate was the dominant factor:
F\textsubscript{1} increased monotonically with LR over the tested
range (5e-6 to 5e-5) for both bases. Layer-wise LR decay did not
improve any cell at any tested value (1.0, 0.95, 0.9). The selected
configuration is therefore LR = 5$\times$10\textsuperscript{$-$5}
without LLRD.

\subsection{Training procedure}

The training mix is the train split of the released dataset. An
earlier build of the model used an AI4Privacy English/email
augmentation slice on top of the in-domain training data to shore
up English-anchor coverage and CH-region email recall; the
released checkpoint trains only on the in-domain split so the
pipeline is end-to-end reproducible from the gheim repository
alone. Hyperparameters are listed in Table~\ref{tab:hyperparams}.

\begin{table}[h]
\centering
\caption{Training hyperparameters.}
\label{tab:hyperparams}
\small
\begin{tabular}{l l}
\toprule
Hyperparameter        & Value \\
\midrule
Optimiser             & AdamW~\cite{loshchilov2019decoupled} \\
Learning rate         & 5e-5 \\
LR scheduler          & Cosine, 5\% warmup \\
Layer-wise LR decay   & 1.0 (off) \\
Effective batch size  & 128 (per-device 64 $\times$ 2 GPUs DDP) \\
Epochs                & 3 max; best at step 3{,}500 of 3{,}987 (epoch 2.63) \\
Max sequence length   & 512 \\
Precision             & bf16 \\
Hardware              & 2 $\times$ NVIDIA RTX 4090 \\
Wall time             & ${\approx}$ 66 min train + 5 min eval \\
\bottomrule
\end{tabular}
\end{table}

The best checkpoint was selected by \texttt{eval\_overall\_f1} on the
validation split. The held-out test split was scored exactly once on
the selected checkpoint to produce the headline numbers in
Section~\ref{sec:eval}.

\subsection{Deployment formats and SDKs}

The model ships in two formats. Server-side users load
\texttt{model.safetensors} (fp32 PyTorch, 2.2 GB) via
\textsf{transformers}~\cite{wolf2020transformers}; in-browser users load
\texttt{onnx/model\_quantized.onnx} (int8 dynamic-quantised ONNX,
552 MB) via \textsf{transformers.js}~\cite{transformersjs} on WebGPU
or WebAssembly. Two thin SDKs wrap the model for the
LLM-anonymisation round-trip use case: \textsf{gheim} on PyPI
(Python) and \textsf{gheim} on npm (TypeScript / JavaScript). Both
expose the same primitives (anonymise text or chat messages,
de-anonymise the streamed response, and a session object that
preserves the sentinel mapping across calls), and both default to
the calibrated decoder discussed in Section~\ref{sec:limits}. A
live in-browser demo at \href{https://gheim.ch}{gheim.ch} runs the
ONNX model through the JS SDK end-to-end, with no server-side
inference. The int8 ONNX export reproduces fp32 strict
F\textsubscript{1} within 0.6~pp on the held-out test split
(0.9044 vs 0.9105) and char F\textsubscript{1} within 0.13~pp
(0.9448 vs 0.9461), at 557~MB for the quantised file versus 2.2~GB
for the fp32 ONNX. Per-category char-F\textsubscript{1} deltas are
all within $\pm 0.003$~pp; the int8 loss is concentrated almost
entirely on \texttt{private\_address} ($-$0.003~pp char), the same
cell that is weakest under fp32. Structured-PII categories
(\texttt{account\_number}, \texttt{private\_email},
\texttt{private\_phone}, \texttt{private\_url},
\texttt{secret}) are unaffected by the int8 step.

% --------------------------------------------------------------------
\section{Evaluation}
\label{sec:eval}

\subsection{Methodology}

Two metrics are reported throughout. \textbf{Strict-span F\textsubscript{1}}
is the standard NER literature metric: a predicted span counts as a
true positive only if its (start, end, label) tuple matches a gold span
exactly. Computed via \textsf{seqeval}~\cite{nakayama2018seqeval} at
the token level for HF and ONNX backends, and via char-level set match
for span-emitting backends (spaCy, Presidio) that lack a token grid.
\textbf{Char-level F\textsubscript{1}} (label-aware) is computed over
$(\text{character index}, \text{label})$ pairs: it asks ``did the model
mask the right characters with the right category?'' and is invariant
to entity-segmentation policy differences across systems.

The latter matters because cross-model PII evaluation has a structural
problem: different annotation conventions split the same surface form
differently. \textsf{ai4privacy/openpii-1m} splits a person mention
into a \texttt{GIVENNAME} and a \texttt{SURNAME} span;
\textsf{gheim-ch-560m} and CoNLL-2003 emit one combined \texttt{PERSON}
span; address detectors fragment ``Werdstrasse 36, 8004 Zürich''
at the comma or do not. Strict-span F\textsubscript{1} penalises every
boundary mismatch even when the underlying redaction-utility goal is
met. Reporting both metrics lets the reader judge which is appropriate
for their use case: strict for literature comparability, char for
deployment.

\subsection{In-distribution performance}

The numbers in this subsection and the next come from a held-out
test split drawn from the same Apertus source corpora as the
training data, with strict document-level disjointness between
splits but the same sampling distribution. They are an
in-distribution upper bound rather than a generalization signal;
the cross-domain numbers in Section~\ref{sec:crossdomain} are the
more meaningful evidence that the model transfers to data drawn
from a different distribution.

On the held-out test split of \textsf{gheim-ch-pii-212k} (21{,}246
chunks), \textsf{gheim-ch-560m} attains a strict-span F\textsubscript{1}
of 0.910 (precision 0.890, recall 0.932) and a char F\textsubscript{1}
of 0.946. The validation/test F\textsubscript{1} gap is 0.000,
consistent with no observable overfitting to the validation set
used for checkpoint selection. The full per-language $\times$
per-category char F\textsubscript{1} breakdown is in
Table~\ref{tab:perlangcat} (Appendix~\ref{app:results}).

Structured categories with strong surface signals
(\texttt{secret}, \texttt{email}, \texttt{phone}, \texttt{URL},
\texttt{account\_number}) hold above 0.97 char F\textsubscript{1}
on average across languages. \texttt{private\_person} and
\texttt{private\_address} carry the per-language variance, with
character-level F\textsubscript{1} ranging 0.85--0.96 across the
five languages. Romansh is the weakest of the five languages
(0.92 char overall, 0.85 char on address) and reflects the smaller
Romansh corpus available for training. For production use the
\texttt{account\_number} and \texttt{private\_address} categories
are intended to be paired with a deterministic regex front-end that
applies checksum validation (see
\texttt{packages\Bsl gheim\Bhy py\Bsl src\Bsl gheim\Bsl detectors\Bsl composite.py}).

\subsection{Comparison to other systems}

Eight other open PII / NER systems were scored on the same test split
under the same harness; the full result table is in
Table~\ref{tab:directionA} (Appendix~\ref{app:results}).
PER-only models (Davlan, dslim, spaCy) emit only person mentions and
no other PII categories, so their overall 8-category
F\textsubscript{1} is reported as n/a; the meaningful number for
those is the \texttt{private\_person} cell.
\textsf{openai/privacy-filter} is scored via its shipped ONNX export
because its custom \texttt{model\_type=openai\_privacy\_filter} is
not registered in \textsf{transformers} 4.57; the harness patches
\texttt{config.json}'s model\_type to \texttt{xlm-roberta} so the
auto-config loader can read the id-to-label mapping, then runs
inference through the unmodified ONNX graph.
\textsf{ZurichNLP/swissbert-ner} is run via an XMOD-adapter-aware
wrapper (the model dispatches per chunk's language) plus a regex
fallback for the non-NER categories (IBAN, AHV, VAT-CHE, phone,
email, URL, secret).

Of the public detectors we measured on Swiss text in this
configuration, both released \textsf{gheim-ch-560m} checkpoints are
the strongest on every cell that is directly comparable. The
comparison is asymmetric: the \textsf{gheim} checkpoints are
in-distribution on this test set (the training data was drawn from
the same corpora) and the contestants are all run zero-shot, so the
absolute gap is partly a domain-fit effect, not only a model-quality
difference. With that caveat noted, the two metrics tell the same
ranking story: among the contestants the strongest
\texttt{private\_person} char F\textsubscript{1} comes from
\textsf{ZurichNLP/swissbert-ner} + regex (0.770) and
\textsf{dslim/bert-base-NER}~\cite{dslim2020bertbasener}
(0.834, English slice only); under
strict F\textsubscript{1} \textsf{openai/privacy-filter} leads at
0.449. The gap between strict and char is largest for systems whose
schema differs most from ours:
\textsf{Isotonic/distilbert\_finetuned\_ai4privacy\_v2} gains
$+30$~pp under char F\textsubscript{1} because its fine-grained
schema (FIRSTNAME, LASTNAME, MIDDLENAME, PREFIX, $\ldots$) collapses
under our 8-category remap and produces highly fragmented spans.

Among the contestants, \textsf{gheim-ch-560m} is also the only entry
with native Romansh support; spaCy has no Romansh or English
pipelines in this configuration, and \textsf{dslim/bert-base-NER} is
English-only. Section~\ref{sec:crossdomain} reports the same model
on each contestant's training distribution for the reverse
comparison.

\subsection{Cross-domain evaluation}
\label{sec:crossdomain}

The released model was scored on four established external benchmarks
(\textsf{ai4privacy/openpii-1m}, Babelscape \textsf{WikiNeural},
ZurichNLP \textsf{swissner}, CoNLL-2003); the full result table is
in Table~\ref{tab:directionB}
(Appendix~\ref{app:results}). For pure-NER datasets
(\textsf{swissner}, CoNLL-2003, WikiNeural) only PER maps to a
\textsf{gheim} category, so the meaningful number is the PER cell.
To anchor the gheim cells, the table reports alongside each row the
F\textsubscript{1} of a model originally trained on that benchmark's
distribution (reproduced through our harness).

Row by row, \textsf{gheim-ch-560m} transfers well to three of the
four benchmarks despite being run zero-shot and without ever seeing
training data from any of them: PER character-F\textsubscript{1}
is 0.911 on CoNLL-2003, 0.808 on Babelscape WikiNeural, and 0.938
on AI4Privacy openpii-1m. The notable strict-vs.-char gap on
\textsf{openpii-1m} (0.178 PER strict, 0.938 PER char) is a schema
artefact: openpii-1m splits each person into separate
\texttt{GIVENNAME} and \texttt{SURNAME} gold spans, which fragment
into two adjacent \texttt{private\_person} spans under our remap
even though every PII character is correctly masked.

The cross-domain row that the released \textsf{gheim-ch-560m}
checkpoint does \emph{not} carry strongly is Swiss news on
\textsf{swissner}: per-language person-character F\textsubscript{1}
is 0.54 (de), 0.76 (fr), 0.64 (it), 0.41 (rm). The released
checkpoint trains only on the in-domain
\textsf{gheim-ch-pii-212k} corpus so that the published pipeline
is end-to-end reproducible from this repository alone; the
swissner gap is the price of that choice. We ran a controlled
multi-source training experiment to quantify the gap and to test
whether a commercial-licence-only training mix could close it.
Section~\ref{sec:multisource} reports the three-variant comparison
and the two-checkpoint release decision that follows from it.

\subsection{Multi-source training experiment and the two-checkpoint release}
\label{sec:multisource}

The architecture is held fixed (\textsf{xlm-roberta-large}, 560M
parameters, same hyperparameters as Table~\ref{tab:hyperparams})
and only the training mix is varied. Three variants are compared:

\begin{itemize}
  \item \textbf{baseline} ($n=170{,}001$): the in-domain train split
        of \textsf{gheim-ch-pii-212k} only.
  \item \textbf{commercial} ($n=239{,}001$): baseline + the train
        splits of \textsf{ai4privacy/openpii-1m} (Apache 2.0, 40k
        chunks in de/fr/it/en) and \textsf{unimelb-nlp/wikiann}
        (CC BY 4.0 / ODC-BY over Wikipedia, ${\approx}$29k chunks in
        de/fr/it/en/rm). All training data is commercially
        relicensable so the resulting model can ship as Apache 2.0.
  \item \textbf{research} ($n=253{,}597$): baseline + the train
        splits of \textsf{ai4privacy/openpii-1m} (40k) +
        Babelscape \textsf{WikiNeural} (CC BY-NC-SA 4.0,
        ${\approx}$39k chunks) + CoNLL-2003 (Reuters
        research-only, 4.4k). The most restrictive upstream
        licence binds the released checkpoint, so this variant is
        non-commercial / research-only.
\end{itemize}

Validation and test splits are held constant across all three
runs (the in-domain val/test of \textsf{gheim-ch-pii-212k}) so
the in-distribution numbers are directly comparable. External
benchmarks are scored zero-shot on each checkpoint. The full
result is in Table~\ref{tab:multisource} (Appendix~\ref{app:multisource}).

In-distribution, the three checkpoints are statistically
indistinguishable (strict-span F\textsubscript{1} = 0.910 / 0.912
/ 0.912; char F\textsubscript{1} = 0.946 across all three).
Out-of-distribution, the picture is asymmetric:

\begin{itemize}
  \item On \textsf{swissner} (Swiss-news NER, zero-shot for all
        three variants — there is no swissner train split), the
        commercial mix lifts \texttt{private\_person} char
        F\textsubscript{1} from 0.70 (baseline) to 0.84 (+14~pp);
        research lifts it to 0.90 (+20~pp). The commercial gain
        is real but recovers only about half of the research
        gain.
  \item On \textsf{openpii-1m} and \textsf{open-pii-500k} (PII
        benchmarks that share schema family with our label space),
        both commercial and research lift character-level
        F\textsubscript{1} into the 0.95--1.0 range; baseline
        stays at 0.87--0.90.
  \item On CoNLL-2003 and WikiNeural (news / Wikipedia NER), both
        the commercial and the research variants \emph{regress}
        against the in-domain-only baseline on the
        \texttt{private\_person} cell: the commercial mix loses
        $-$36~pp on CoNLL and $-$13~pp on WikiNeural; the research
        mix loses $-$15~pp on CoNLL and $-$1~pp on WikiNeural, even
        though it trains on the train splits of both. The cause is
        the same in both cases — the broader 8-category output
        schema of the gheim checkpoints produces non-PER false
        positives on news / Wikipedia text (where any IBAN-shaped,
        email-shaped, or URL-shaped string is flagged) that the
        narrower in-domain-only baseline does not produce. The
        result is a small but real precision tradeoff for the
        wider category coverage these variants target.
  \item On \textsf{gretelai/synthetic\_pii\_finance\_multilingual}
        (multilingual financial-document PII, Apache 2.0), all three
        checkpoints converge near 0.60 character-F\textsubscript{1}.
        No training mix contains financial-document text; the
        benchmark therefore measures the architecture's
        structural-PII transfer rather than mix-induced
        differences. The flat result motivates adding this dataset
        to a future commercial training mix.
\end{itemize}

The conclusion of the experiment is that no single training mix
strictly dominates: the commercial mix improves swissner at the
cost of CoNLL-2003 / WikiNeural; the research mix improves
swissner and openpii cleanly but is non-commercial; the in-domain
baseline ships under the cleanest licence but loses on Swiss-news
transfer. We therefore publish two checkpoints from this paper:

\begin{itemize}
  \item \textsf{joelbarmettler/gheim-ch-560m} — the in-domain
        baseline, released under \textbf{Apache 2.0} for production
        and commercial deployments. Cross-domain swissner PER
        char F\textsubscript{1} 0.70.
  \item \textsf{joelbarmettler/gheim-ch-560m-research} — the
        multi-source research variant, released under
        \textbf{CC BY-NC-SA 4.0} with a Reuters research-only
        rider (from CoNLL-2003), for academic research /
        benchmarks / non-commercial evaluation. Cross-domain
        swissner PER char F\textsubscript{1} 0.89.
\end{itemize}

The commercial-mix variant is not published as a separate
checkpoint --- it does not strictly dominate the baseline on
cross-domain, and the half-recovery on swissner does not justify
a third release. Closing the commercial swissner gap further
will require adding more news-domain, commercially-relicensable
NER training data (candidates: \textsf{MultiCoNER v2} under
CC BY 4.0, GermEval 2014 under CC BY-SA 4.0,
\textsf{gretelai/synthetic\_pii\_finance\_multilingual} under
Apache 2.0); we leave that to a future iteration.

A methodology check that replicates each baseline on its own
training-distribution dataset and compares the reproduction to its
published number is reported in Appendix~\ref{app:results}
(Table~\ref{tab:replication}). Three of four reproductions match
within $\pm$0.03; the fourth (Isotonic) is a metric mismatch
(per-token accuracy versus strict-span F\textsubscript{1}) rather
than a harness bug.

% --------------------------------------------------------------------
\section{Limitations}
\label{sec:limits}

\begin{itemize}
  \item \textbf{Recall-oriented labelling policy.} Public Swiss documents
        contain extensive personally-identifiable information that is
        already public (court rulings, parliament records, public
        officials). The dataset labels these as PII and the model
        learns to flag them. Applications that need to distinguish
        public from private entities should layer a downstream
        public-vs-private classifier.
  \item \textbf{\texttt{private\_address} test strict F\textsubscript{1}
        is 0.84} (char F\textsubscript{1} 0.92). Boundary placement
        on multi-token addresses is the dominant error mode.
  \item \textbf{\texttt{account\_number} test strict F\textsubscript{1}
        is 0.99} in the headline, but a small fraction of regex-shaped
        non-PII (numeric tables in court documents and parliamentary
        statistics) still slips through. Production deployments
        should pair the model with a regex front-end with checksum
        validation (IBAN, AHV, VAT-CHE, Luhn).
  \item \textbf{Romansh test strict F\textsubscript{1} is 0.89}
        (char F\textsubscript{1} 0.93), the weakest of the five
        languages. The Romansh training material is dominated by a
        single literary/journalistic register (Rumantsch Grischun),
        so dialectal or technical RM text is unmeasured.
  \item \textbf{Cross-domain Swiss-news transfer is meaningfully
        weaker than the in-distribution headline}, with
        \textsf{swissner} per-language person-character
        F\textsubscript{1} dropping to 0.54 (de) / 0.76 (fr) /
        0.64 (it) / 0.41 (rm). The released checkpoint trains
        only on the in-domain corpus; an alternative training mix
        that includes external public NER corpora is the planned
        next iteration.
  \item \textbf{Swiss German dialect (GSW) is not measured.} The
        fasttext detector used in data preparation labels GSW as
        standard German.
  \item \textbf{Annotations are machine-generated.} The dataset
        was not human-annotated; the four-LLM OpenRouter audit in
        Section~\ref{sec:labelnoise} (F\textsubscript{1}~0.71
        overall against the majority-of-four consensus, with
        policy-interpretation disagreements dominating the gap)
        is the only quantitative labelling-policy reference we
        can quote.
  \item \textbf{Re-identification is not in scope.} The model is
        intended for redaction. It is not designed to identify or
        link entities to real persons and does not return
        entity-linked identifiers.
\end{itemize}

\paragraph{Inference-time suppression and the calibrated decoder.}
A failure mode visible only on deploy probes, not in the test-set
metrics: on a small fraction of inputs the raw model emits zero
non-O tokens for an entire chunk, even when the chunk contains
clear PII. The pattern reproduces consistently when a low-confidence
person mention is immediately followed by a structurally distinctive
token (e.g.\ a CH-IBAN), which appears to nudge the per-token
argmax over the entire span toward the dominant \texttt{O} class.
The released SDKs ship a thin wrapper that subtracts a small
constant bias from the \texttt{O}-class logit before argmax
(\texttt{CalibratedDetector}, default bias 0.5). On the failing
probe set this recovers approximately 95\% of the suppression
cases; on the in-distribution test split it costs approximately
0.4 pp strict F\textsubscript{1}. The wrapper is the default
detector in both the Python and JavaScript SDKs. The published
checkpoint exhibits a known inference-time failure mode that the
test-set numbers do not expose; the wrapper trades a small amount
of in-distribution F\textsubscript{1} for a noticeable improvement
on the failure mode without requiring retraining.

% --------------------------------------------------------------------
\section{Summary}
\label{sec:conclusion}

This report described \textsf{gheim-ch-pii-212k} and
\textsf{gheim-ch-560m}, a Swiss-tuned PII NER dataset and model
targeting LLM-API anonymisation round-tripping, a deployment
context that existing open detectors handle poorly on Swiss text.
On the in-distribution test split both released checkpoints attain
a strict-span F\textsubscript{1} of 0.91 (char F\textsubscript{1}
0.95); among the public detectors we measured zero-shot on the
same set the strongest overall is Microsoft Presidio at 0.398
strict / 0.501 char. On out-of-distribution Swiss news
(\textsf{ZurichNLP/swissner}, the only Swiss-news NER benchmark that
includes Romansh), the Apache-2.0 baseline attains
\texttt{private\_person} character F\textsubscript{1} of 0.70 and
the research variant attains 0.90 — a 20~pp gain from broader
training data even though swissner itself is not used at training
time (no swissner train split exists; the gain comes from
generalisation off the additional openpii-1m / WikiNeural /
CoNLL-2003 supervision in the research mix). The dataset is released under
CC BY 4.0; the model under Apache 2.0; both are on the Hugging Face
Hub. The full evaluation harness, the per-system JSON outputs, and
the reproducible build of the comparison tables are at
\url{https://github.com/joelbarmettlerUZH/gheim}.

\paragraph{Known follow-on work.} The dataset has no coverage of
spoken Swiss German (the fasttext language identifier labels GSW as
standard German) and only one Romansh register (Rumantsch Grischun);
both gaps would require additional source corpora rather than
pipeline changes. The model's released checkpoint is 560M
parameters, larger than is convenient for edge or browser
deployment; a smaller distilled variant would be useful but is
not part of this release. The dataset's recall-oriented labelling
policy flags publicly-listed institutional entities (e.g.\ a court
switchboard phone) as PII; applications that require stricter
precision would benefit from a downstream public-vs-private entity
classifier layered on top. The label-noise audit in
Section~\ref{sec:labelnoise} also pointed at three concrete
dataset-cleanup directions for a future revision: (a) tightening
the regex-layer \texttt{account\_number} column-context heuristic
to catch the residual false positives in numeric tables, (b)
optionally exposing the recall-aggressive policy on court-document
roles as an opt-in flag rather than the default, and (c) further
hardening the commercial-register re-labelling pass so director
and board-member names stop being under-recalled in long
contact-block listings.

\bibliographystyle{ACM-Reference-Format}
\bibliography{references}

\appendix

\section{Detailed evaluation results}
\label{app:results}

This appendix contains the four full result tables referenced from
Section~\ref{sec:eval}: the in-distribution per-language $\times$
per-category breakdown for \textsf{gheim-ch-560m}; the comparison to
seven other public detectors on the same test split; the cross-domain
transfer numbers on four external benchmarks; and a methodology
check that re-runs each external baseline on its own training
distribution. The prose in Section~\ref{sec:eval} summarises what
these tables show; the tables themselves are kept out of the main
narrative so each can be inspected at full width.

\subsection{In-distribution per-language $\times$ per-category}
\label{app:perlangcat}

Table~\ref{tab:perlangcat} gives the per-language $\times$ per-category
char-level F\textsubscript{1} on the held-out test split of
\textsf{gheim-ch-pii-212k}. Structured categories with strong surface
signals (\texttt{secret}, \texttt{email}, \texttt{phone},
\texttt{URL}) hold above 0.95 in every language;
\texttt{account\_number} carries most of the per-language variance,
with the Romansh cell as the only one below 0.5.

\begin{table*}[!htbp]
\centering
\caption{Per-language $\times$ per-category char-level
F\textsubscript{1} on the held-out test split of
\textsf{gheim-ch-pii-212k} under \textsf{gheim-ch-560m}.
The right-most column gives the per-category average
(over languages, gold-weighted); the bottom row gives the
per-language average (over categories).}
\label{tab:perlangcat}
\small
\begin{tabular}{l r r r r r @{\hspace{1em}} r}
\toprule
Category                    & de\_ch & fr\_ch & it\_ch & rm    & en    & Avg.\ \\
\midrule
\texttt{account\_number}    & 0.994  & 0.998  & 0.990  & 0.971 & 1.000 & 0.994 \\
\texttt{private\_address}   & 0.933  & 0.911  & 0.916  & 0.853 & 0.955 & 0.917 \\
\texttt{private\_date}      & 0.951  & 0.908  & 0.952  & 0.919 & 0.883 & 0.933 \\
\texttt{private\_email}     & 0.996  & 1.000  & 0.997  & 0.989 & 1.000 & 0.997 \\
\texttt{private\_person}    & 0.913  & 0.939  & 0.951  & 0.909 & 0.955 & 0.930 \\
\texttt{private\_phone}     & 0.995  & 0.996  & 0.996  & 0.985 & 1.000 & 0.995 \\
\texttt{private\_url}       & 0.990  & 0.993  & 0.993  & 0.994 & 0.970 & 0.991 \\
\texttt{secret}             & 0.994  & 0.999  & 1.000  & 0.999 & n/a   & 0.997 \\
\midrule
Avg.\                       & 0.944  & 0.931  & 0.956  & 0.918 & 0.954 & 0.940 \\
\bottomrule
\end{tabular}
\end{table*}

\subsection{Comparison to other detectors}
\label{app:directionA}

Table~\ref{tab:directionA} reports seven other open PII / NER systems
on the \textsf{gheim-ch-pii-212k} test split under the same harness.
Both the overall 8-category metric and the \texttt{private\_person}
cell are shown for each system; PER-only systems have n/a for the
overall metric since they cannot emit the other seven categories.
The comparison is in-distribution for \textsf{gheim-ch-560m} and
zero-shot for the seven contestants; see
Section~\ref{sec:crossdomain} for the reverse direction.

\begin{table*}[!htbp]
\centering
\caption{Other models on the \textsf{gheim-ch-pii-212k} test split.
Strict and char F\textsubscript{1} for both the overall 8-category
metric and the \texttt{private\_person} cell.}
\label{tab:directionA}
\small
\begin{tabular}{l r r r r}
\toprule
Model                                                            & Strict F1 & Char F1 & PER strict & PER char \\
\midrule
\textbf{\textsf{joelbarmettler/gheim-ch-560m}}                  & \textbf{0.910} & \textbf{0.946} & \textbf{0.896} & \textbf{0.929} \\
\textbf{\textsf{joelbarmettler/gheim-ch-560m-research}}          & \textbf{0.912} & \textbf{0.946} & \textbf{0.898} & \textbf{0.929} \\
\textsf{dslim/bert-base-NER} (English slice)                     & n/a   & n/a   & 0.561 & 0.834 \\
\textsf{ZurichNLP/swissbert-ner} + regex hybrid                  & 0.346 & 0.613 & 0.464 & 0.770 \\
\textsf{Davlan/distilbert-base-multilingual-cased-ner-hrl}       & n/a   & n/a   & 0.569 & 0.758 \\
\textsf{Davlan/xlm-roberta-base-ner-hrl}                         & n/a   & n/a   & 0.610 & 0.733 \\
\textsf{openai/privacy-filter} (1.4B MoE, ONNX, zero-shot)       & 0.449 & 0.619 & 0.436 & 0.582 \\
spaCy \texttt{de/fr/it\_core\_news\_lg}                          & n/a   & n/a   & 0.473 & 0.574 \\
Microsoft Presidio (multi-lang config)                           & 0.398 & 0.501 & 0.433 & 0.519 \\
\textsf{Isotonic/distilbert\_finetuned\_ai4privacy\_v2}          & 0.140 & 0.441 & 0.243 & 0.513 \\
\bottomrule
\end{tabular}
\end{table*}

\subsection{Cross-domain transfer}
\label{app:directionB}

Table~\ref{tab:directionB} reports \textsf{gheim-ch-560m} PER
F\textsubscript{1} on four external benchmarks the model has never
seen. To anchor each row, the same column reports a model trained
on (or claimed-state-of-the-art on) that benchmark's native
distribution, reproduced through our harness. The footnotes give
the specific metric / slice each native number is computed on.

\begin{table*}[!htbp]
\centering
\caption{\textsf{gheim-ch-560m} PER F\textsubscript{1} on four
external benchmarks, anchored against a specialist model trained or
claimed-state-of-the-art on each benchmark's native distribution.
\textsf{gheim-ch-560m} is run zero-shot; native baselines are
reproduced through the same harness.}
\label{tab:directionB}
\small
\begin{tabular}{l r r r l r}
\toprule
External benchmark                                & $n$       & gheim PER strict & gheim PER char & Native baseline                              & Native F\textsubscript{1} \\
\midrule
\textsf{ai4privacy/pii-masking-openpii-1m}        & 8{,}000   & 0.178            & 0.938          & \textsf{Isotonic\Bsl distilbert\_ai4privacy} & 0.78$^\dagger$  \\
\textsf{ai4privacy/open-pii-masking-500k}         & 8{,}000   & 0.540            & 0.933          & --- (sibling of openpii-1m)                  & ---             \\
\textsf{gretelai/synthetic\_pii\_finance\_multilingual} & 4{,}800 & 0.519          & 0.624          & --- (no public baseline)                     & ---             \\
Babelscape \textsf{WikiNeural}                    & 8{,}000   & 0.760            & 0.808          & \textsf{Davlan\Bsl xlm\Bhy r\Bhy base\Bhy ner\Bhy hrl} & 0.813$^\ddagger$ \\
ZurichNLP \textsf{swissner}                       & 800       & 0.637            & 0.702          & spaCy \texttt{de\_core\_news\_lg}            & 0.841$^\S$      \\
CoNLL-2003                                        & 3{,}453   & 0.854            & 0.911          & \textsf{dslim/bert-base-NER}~\cite{dslim2020bertbasener}      & 0.957           \\
\bottomrule
\end{tabular}
\end{table*}

\noindent$^\dagger$Isotonic's strict-span F\textsubscript{1} on its own
training distribution; per-token accuracy across all categories
(the metric Isotonic itself reports) is 0.969 in our reproduction
versus 0.955 claimed.
$^\ddagger$Davlan-XLMR's overall multilingual avg from our
reproduction (claimed ${\sim}$0.84).
$^\S$spaCy de\_core's overall F\textsubscript{1} on the German slice
of \textsf{swissner} (spaCy ships per-language pipelines and has no
single multilingual model that covers the whole set).

\clearpage
\subsection{Multi-source training experiment}
\label{app:multisource}

Table~\ref{tab:multisource} reports the full three-variant
comparison summarised in Section~\ref{sec:multisource}. In-domain
numbers are character-level F\textsubscript{1} on the held-out
test split of \textsf{gheim-ch-pii-212k}; cross-domain numbers are
PER character-level F\textsubscript{1} on each external benchmark.

\begin{table*}[!htbp]
\centering
\caption{Three training-mix variants compared. ``Baseline''
trains on in-domain data only and ships as
\textsf{gheim-ch-560m} (Apache 2.0). ``Research'' adds three
external NER / PII corpora and ships as
\textsf{gheim-ch-560m-research} (CC BY-NC-SA 4.0 with a Reuters
research-only rider). ``Commercial'' (openpii + WikiAnn only)
is reported as an experimental data point and is not separately
released.}
\label{tab:multisource}
\small
\begin{tabular}{l r r r}
\toprule
                                                & Baseline (Apache 2.0) & Commercial (experimental) & Research (CC BY-NC-SA) \\
\midrule
Train chunks                                    & 170{,}001 & 239{,}001 & 253{,}597 \\
External training corpora                       & ---       & openpii + WikiAnn & openpii + WikiNeural + CoNLL-2003 \\
\midrule
In-domain test strict F\textsubscript{1}        & 0.910     & 0.912     & 0.912     \\
In-domain test char F\textsubscript{1}          & 0.946     & 0.946     & 0.946     \\
\midrule
\multicolumn{4}{l}{\emph{\texttt{private\_person}-cell character-level F\textsubscript{1} on each external benchmark}} \\
\textsf{ZurichNLP/swissner} (zero-shot)          & 0.702     & 0.844     & \textbf{0.903} \\
\quad de                                        & 0.539     & 0.805     & \textbf{0.931} \\
\quad fr                                        & 0.761     & 0.845     & \textbf{0.913} \\
\quad it                                        & 0.643     & 0.745     & \textbf{0.856} \\
\quad rm                                        & 0.409     & 0.552     & \textbf{0.873} \\
\textsf{ai4privacy/openpii-1m}                  & 0.938     & 0.995     & 0.995     \\
\textsf{ai4privacy/open-pii-500k}               & 0.933     & 0.973     & 0.982     \\
\textsf{gretelai/synthetic\_pii\_finance}       & 0.624     & 0.619     & 0.627     \\
Babelscape \textsf{WikiNeural}                  & 0.808     & 0.681$^\downarrow$ & 0.795$^\downarrow$ \\
CoNLL-2003                                      & 0.911     & 0.548$^\downarrow$ & 0.765$^\downarrow$ \\
\bottomrule
\end{tabular}
\end{table*}

\noindent$^\downarrow$The commercial variant regresses against
baseline on these two benchmarks because the substitute
\textsf{unimelb-nlp/wikiann} training data is short Wikipedia
spans, lower per-chunk signal density than
\textsf{WikiNeural}-train or CoNLL-train; the research variant
includes both as training material and so does not regress.

\subsection{Methodology validation: replicating each baseline on its own training distribution}
\label{app:replication}

To check that the harness in Section~\ref{sec:eval} is sound rather
than an unfair implementation, we replicated each external system on
the dataset on which it was trained and compared our reproduction to
the system's own published F\textsubscript{1}
(Table~\ref{tab:replication}). Three of four reproductions match
within $\pm$0.03; the fourth (Isotonic) is a metric mismatch
discussed below.

\begin{table*}[!htbp]
\centering
\caption{Each baseline scored on its own training-distribution
dataset, against published numbers.}
\label{tab:replication}
\small
\begin{tabular}{l r r r}
\toprule
Replication target                                & Claimed       & Repro             & $\Delta$ \\
\midrule
dslim/bert-base-NER on CoNLL-2003                 & 0.910         & 0.913             & $+0.003$ \\
dslim PER on CoNLL-2003                           & ${\sim}$0.96  & 0.957             & $-0.003$ \\
Davlan-XLMR on WikiNeural avg                     & ${\sim}$0.84  & 0.813             & $-0.027$ \\
spaCy de\_core on swissner\_de                    & ${\sim}$0.83  & 0.841             & $+0.011$ \\
Isotonic distilbert on AI4Privacy                 & 0.955         & 0.969$^\dagger$   & $+0.014$ \\
\bottomrule
\end{tabular}
\end{table*}

\noindent$^\dagger$Isotonic's claimed
F\textsubscript{1} of 0.955 initially reproduced as 0.78 strict-span
F\textsubscript{1} on the same data, an apparent 17-pp gap. The
difference traces to the metric: Isotonic's reported number is
\emph{per-token classification accuracy} (which includes the dominant
\texttt{O} class and is much higher than per-span F\textsubscript{1}).
Reproducing Isotonic's metric exactly (per-token accuracy including
\texttt{O}) yields 0.969, matching their 0.955 within sampling noise.
Our strict-span F\textsubscript{1} of 0.78 on the same data is the
correct NER-literature comparison.

\end{document}
